\documentclass[10pt,a4paper]{amsart}
\usepackage{amsmath,amssymb,amsthm,graphicx}
% Retained preamble for the Stacks Project chapter snapshot in this directory.
% The chapter's own preamble.tex is a copy of the shared Stacks Project preamble; the block
% below is sliced MECHANICALLY, line for line, from it. Nothing was added, removed or
% redefined; the wrapper class, the xr-hyper/hyperref wiring and the externaldocument list
% are omitted because the assembler supplies the document class.
%
% lines taken from preamble.tex: 3, 5, 9, 11, 12, 156, 157, 158, ... (51 lines total)
%
% package coverage restored from the chapter's own preamble (2026-10-13): xy, xy-2cell

\usepackage[all]{xy}\usepackage{verbatim}
\xyoption{2cell}
\UseAllTwocells

\newenvironment{reference}{\comment}{\endcomment}
\newenvironment{slogan}{\comment}{\endcomment}
\newenvironment{history}{\comment}{\endcomment}
\theoremstyle{plain}
\newtheorem{theorem}[subsection]{Theorem}
\newtheorem{proposition}[subsection]{Proposition}
\newtheorem{lemma}[subsection]{Lemma}
\theoremstyle{definition}
\newtheorem{definition}[subsection]{Definition}
\newtheorem{example}[subsection]{Example}
\newtheorem{exercise}[subsection]{Exercise}
\newtheorem{situation}[subsection]{Situation}
\theoremstyle{remark}
\newtheorem{remark}[subsection]{Remark}
\newtheorem{remarks}[subsection]{Remarks}
\numberwithin{equation}{subsection}
\def\lim{\mathop{\mathrm{lim}}\nolimits}
\def\colim{\mathop{\mathrm{colim}}\nolimits}
\def\Spec{\mathop{\mathrm{Spec}}}
\def\Hom{\mathop{\mathrm{Hom}}\nolimits}
\def\Ext{\mathop{\mathrm{Ext}}\nolimits}
\def\SheafHom{\mathop{\mathcal{H}\!\mathit{om}}\nolimits}
\def\SheafExt{\mathop{\mathcal{E}\!\mathit{xt}}\nolimits}
\def\Sch{\mathit{Sch}}
\def\Mor{\mathop{\mathrm{Mor}}\nolimits}
\def\Ob{\mathop{\mathrm{Ob}}\nolimits}
\def\Sh{\mathop{\mathit{Sh}}\nolimits}
\def\NL{\mathop{N\!L}\nolimits}
\def\CH{\mathop{\mathrm{CH}}\nolimits}
\def\proetale{{pro\text{-}\acute{e}tale}}
\def\etale{{\acute{e}tale}}
\def\QCoh{\mathit{QCoh}}
\def\Ker{\mathop{\mathrm{Ker}}}
\def\Im{\mathop{\mathrm{Im}}}
\def\Coker{\mathop{\mathrm{Coker}}}
\def\Coim{\mathop{\mathrm{Coim}}}
\DeclareMathSymbol{\boxtimes}{\mathbin}{AMSa}{"02}
\def\QCohstack{\mathcal{QC}\!\mathit{oh}}
\def\Cohstack{\mathcal{C}\!\mathit{oh}}
\def\Spacesstack{\mathcal{S}\!\mathit{paces}}
\def\Quotfunctor{\mathrm{Quot}}
\def\Hilbfunctor{\mathrm{Hilb}}
\def\Curvesstack{\mathcal{C}\!\mathit{urves}}
\def\Polarizedstack{\mathcal{P}\!\mathit{olarized}}
\def\Complexesstack{\mathcal{C}\!\mathit{omplexes}}
\def\Pic{\mathop{\mathrm{Pic}}\nolimits}
\def\Picardstack{\mathcal{P}\!\mathit{ic}}
\def\Picardfunctor{\mathrm{Pic}}
\def\Deformationcategory{\mathcal{D}\!\mathit{ef}}
\usepackage{cleveref}
\crefname{originallemma}{original lemma}{original lemmas}
\Crefname{originallemma}{Original lemma}{Original lemmas}
\makeatletter
\newlabel{derived-lemma-negative-vanishing}{{05TC}{source}{}{sourcepointer-1}{}}
\newlabel{derived-lemma-negative-vanishing@cref}{{[originallemma][0][]05TC}{[1][0][]source}{}{}{}}
\makeatother
\crefname{originallemma}{original lemma}{original lemmas}
\Crefname{originallemma}{Original lemma}{Original lemmas}
\makeatletter
\newlabel{derived-remark-truncation-distinguished-triangle}{{08J5}{source}{}{sourcepointer-2}{}}
\newlabel{derived-remark-truncation-distinguished-triangle@cref}{{[originallemma][0][]08J5}{[1][0][]source}{}{}{}}
\makeatother
\crefname{originallemma}{original lemma}{original lemmas}
\Crefname{originallemma}{Original lemma}{Original lemmas}
\makeatletter
\newlabel{injectives-lemma-RF-commutes-with-Rlim}{{08U1}{source}{}{sourcepointer-3}{}}
\newlabel{injectives-lemma-RF-commutes-with-Rlim@cref}{{[originallemma][0][]08U1}{[1][0][]source}{}{}{}}
\makeatother
\crefname{originallemma}{original lemma}{original lemmas}
\Crefname{originallemma}{Original lemma}{Original lemmas}
\makeatletter
\newlabel{lemma-nice-K-injective}{{08F6}{source}{}{sourcepointer-4}{}}
\newlabel{lemma-nice-K-injective@cref}{{[originallemma][0][]08F6}{[1][0][]source}{}{}{}}
\makeatother
\crefname{originallemma}{original lemma}{original lemmas}
\Crefname{originallemma}{Original lemma}{Original lemmas}
\makeatletter
\newlabel{more-algebra-remark-compare-derived-limit}{{08U5}{source}{}{sourcepointer-5}{}}
\newlabel{more-algebra-remark-compare-derived-limit@cref}{{[originallemma][0][]08U5}{[1][0][]source}{}{}{}}
\makeatother
\crefname{originallemma}{original lemma}{original lemmas}
\Crefname{originallemma}{Original lemma}{Original lemmas}
\makeatletter
\newlabel{more-algebra-section-Rlim}{{07KV}{source}{}{sourcepointer-6}{}}
\newlabel{more-algebra-section-Rlim@cref}{{[originallemma][0][]07KV}{[1][0][]source}{}{}{}}
\makeatother
\title{Stacks, Derived Categories of Spaces: application-nice-K-injective}
\author{The Stacks Project Authors}
\date{stacks-project:spaces-perfect.tex:local-snapshot-20261012; GNU FDL 1.2 or later}
\begin{document}
\maketitle
\noindent\textit{Editorial note.} Statement span is lines 696-717 of sources/stacks-spaces-perfect/pinned-original-source.tex; proof span is lines 719-771, the complete author proof (51 body lines).\par\medskip
\noindent\textit{Editorial note.} Navigation stubs: 6 cited-but-not-included labels are supplied as non-load-bearing pointers with their pinned tags where the snapshot tag index has them.\par\medskip
\noindent\textit{Editorial note.} extra\_packages is empty on purpose: the retained preamble already loads xy with options and a second declaration clashes in the assembler wrapper.\par\medskip
\noindent\textit{Editorial note.} Text is retained verbatim, including the author's own typos and placeholders.\par\medskip
\noindent\textit{Editorial note.} Queue index 1 of queue\_spaceperfect.json (pb=1807, stmt 696-717, proof 719-771).\par\medskip
\noindent\textit{Editorial note.} stacks-spaces-perfect wave sp01 (2026-10-14): fresh adjudication accept per lane-s3/sp01\_judgments; pool per Chief 12:4x unfreeze (reserve-pool workflow, six-check clean); spaces-perfect.tex 126 with-proof rows, 70 b107-occupied labels avoided, 56 free rows judged 32 accept / 24 reject.\par\medskip
\noindent\textit{Editorial note.} Cross-chapter same-label warning (NOT an occupancy conflict, Chief ruling 2026-10-13 evening): label 'lemma-application-nice-K-injective' also exists in stacks-perfect — a different proposition in a different chapter file.\par\medskip
\noindent\textit{Allowed standard prerequisites (editorial).}\begin{itemize}
\item Explicit prerequisite boundary. The cited inputs are prerequisites; the argument in the retained span is not replaced by any cited result.
\end{itemize}\medskip
\ifdefined\hypertarget\hypertarget{sourcepointer-1}{}\fi
\noindent\textit{Source reference.} Original lemma 05TC: sources/stacks-spaces-perfect/pinned-original-source.tex, source lines 771--771. Cited by the retained statement or proof but not part of the collected spans; supplied only as a navigation pointer. The argument does not depend on its content.\par\medskip
\ifdefined\hypertarget\hypertarget{sourcepointer-2}{}\fi
\noindent\textit{Source reference.} Original lemma 08J5: sources/stacks-spaces-perfect/pinned-original-source.tex, source lines 771--771. Cited by the retained statement or proof but not part of the collected spans; supplied only as a navigation pointer. The argument does not depend on its content.\par\medskip
\ifdefined\hypertarget\hypertarget{sourcepointer-3}{}\fi
\noindent\textit{Source reference.} Original lemma 08U1: sources/stacks-spaces-perfect/pinned-original-source.tex, source lines 771--771. Cited by the retained statement or proof but not part of the collected spans; supplied only as a navigation pointer. The argument does not depend on its content.\par\medskip
\ifdefined\hypertarget\hypertarget{sourcepointer-4}{}\fi
\noindent\textit{Source reference.} Original lemma 08F6: sources/stacks-spaces-perfect/pinned-original-source.tex, source lines 669--678. Cited by the retained statement or proof but not part of the collected spans; supplied only as a navigation pointer. The argument does not depend on its content.\par\medskip
\ifdefined\hypertarget\hypertarget{sourcepointer-5}{}\fi
\noindent\textit{Source reference.} Original lemma 08U5: sources/stacks-spaces-perfect/pinned-original-source.tex, source lines 771--771. Cited by the retained statement or proof but not part of the collected spans; supplied only as a navigation pointer. The argument does not depend on its content.\par\medskip
\ifdefined\hypertarget\hypertarget{sourcepointer-6}{}\fi
\noindent\textit{Source reference.} Original lemma 07KV: sources/stacks-spaces-perfect/pinned-original-source.tex, source lines 771--771. Cited by the retained statement or proof but not part of the collected spans; supplied only as a navigation pointer. The argument does not depend on its content.\par\medskip
\begin{lemma}
\label{lemma-application-nice-K-injective}
Let $S$ be a scheme. Let $X$ be an algebraic space over $S$.
Let $F : \textit{Mod}(\mathcal{O}_X) \to \textit{Ab}$
be a functor and $N \geq 0$ an integer. Assume that
\begin{enumerate}
\item $F$ is left exact,
\item $F$ commutes with countable direct products,
\item $R^pF(\mathcal{F}) = 0$ for all $p \geq N$ and $\mathcal{F}$
quasi-coherent.
\end{enumerate}
Then for $E \in D_\QCoh(\mathcal{O}_X)$
\begin{enumerate}
\item $H^i(RF(\tau_{\leq a}E) \to H^i(RF(E))$ is an isomorphism
for $i \leq a$,
\item $H^i(RF(E)) \to H^i(RF(\tau_{\geq b - N + 1}E))$ is an isomorphism
for $i \geq b$,
\item if $H^i(E) = 0$ for $i \not \in [a, b]$ for some
$-\infty \leq a \leq b \leq \infty$, then $H^i(RF(E)) = 0$
for $i \not \in [a, b + N - 1]$.
\end{enumerate}
\end{lemma}

\begin{proof}
Statement (1) is
Derived Categories, Lemma \ref{derived-lemma-negative-vanishing}.

\medskip\noindent
Proof of statement (2). Write $E_n = \tau_{\geq -n}E$. We have
$E = R\lim E_n$, see Lemma \ref{lemma-nice-K-injective}. Thus
$RF(E) = R\lim RF(E_n)$ in $D(\textit{Ab})$ by Injectives, Lemma
\ref{injectives-lemma-RF-commutes-with-Rlim}. Thus for every $i \in \mathbf{Z}$
we have a short exact sequence
$$
0 \to R^1\lim H^{i - 1}(RF(E_n)) \to H^i(RF(E)) \to \lim H^i(RF(E_n)) \to 0
$$
see More on Algebra, Remark
\ref{more-algebra-remark-compare-derived-limit}.
To prove (2) we will show that the term on the left is zero
and that the term on the right equals $H^i(RF(E_{-b + N - 1})$
for any $b$ with $i \geq b$.

\medskip\noindent
For every $n$ we have a distinguished triangle
$$
H^{-n}(E)[n] \to E_n \to E_{n - 1} \to H^{-n}(E)[n + 1]
$$
(Derived Categories, Remark
\ref{derived-remark-truncation-distinguished-triangle})
in $D(\mathcal{O}_X)$. Since $H^{-n}(E)$ is quasi-coherent we have
$$
H^i(RF(H^{-n}(E)[n])) = R^{i + n}F(H^{-n}(E)) = 0
$$
for $i + n \geq N$ and
$$
H^i(RF(H^{-n}(E)[n + 1])) = R^{i + n + 1}F(H^{-n}(E)) = 0
$$
for $i + n + 1 \geq N$. We conclude that
$$
H^i(RF(E_n)) \to H^i(RF(E_{n - 1}))
$$
is an isomorphism for $n \geq N - i$. Thus the systems $H^i(RF(E_n))$ all
satisfy the ML condition and the $R^1\lim$ term in our short exact
sequence is zero (see discussion in
More on Algebra, Section \ref{more-algebra-section-Rlim}).
Moreover, the system $H^i(RF(E_n))$ is constant starting
with $n = N - i - 1$ as desired.

\medskip\noindent
Proof of (3). Under the assumption on $E$ we have
$\tau_{\leq a - 1}E = 0$ and we get the vanishing
of $H^i(RF(E))$ for $i \leq a - 1$ from (1).
Similarly, we have $\tau_{\geq b + 1}E = 0$ and hence
we get the vanishing of $H^i(RF(E))$ for $i \geq b + n$ from
part (2).
\end{proof}

\end{document}
